\chapter{Создание макета сайта чемпионата}
\label{ch:practice}

В данной главе последовательно описано построение макетов всех пяти
смысловых страниц сайта чемпионата \textbf{PHYGITAL~FC} в десктоп-версии
(\(1440\,\)px), а также общие настройки колоночной сетки и привязок,
после чего показаны адаптивные версии для планшета (\(768\,\)px) и
мобильного телефона (\(360\,\)px).

\section{Главная страница}

Главная страница --- витрина чемпионата. Её задача за один экран дать
посетителю понять: что за событие, когда оно состоится, кто играет и
как купить билет. Структура десктоп-фрейма \texttt{Desktop~1440}
страницы \texttt{Главная~страница}:

\begin{itemize}
  \item \textbf{Шапка} (фрейм \texttt{site-header}, высота \(80\,\)px):
        слева --- логотип \textbf{PHYGITAL~FC}; в центре ---
        горизонтальная навигация (\textit{Главная}, \textit{Команды},
        \textit{Турнир}, \textit{Магазин}); справа --- иконка профиля
        и CTA-кнопка \textit{<<Купить билет>>}.
  \item \textbf{Hero-блок} (фрейм \texttt{hero}, высота \(\sim 560\,\)px):
        крупный заголовок \textit{<<Чемпионат по киберфутболу.
        Новосибирск, 2024>>} в шрифте \texttt{Space~Grotesk}
        (\texttt{72\,pt}, жирное); подзаголовок с датой и местом
        проведения; пара CTA-кнопок \textit{<<Купить билет>>} и
        \textit{<<Программа турнира>>}. Фоном --- фотография стадиона
        с накладкой градиента \texttt{\#0B1024} (\(0.7\) opacity).
\end{itemize}

% — Скриншот: главная страница, Hero-блок —
\begin{figure}[h]
  \centering
  \includegraphics[width=0.85\textwidth]{img/03_home_hero}
  \caption{Главная страница: шапка и Hero-блок чемпионата}
  \label{fig:home_hero}
\end{figure}

Ниже Hero-блока размещаются ещё две секции, рассчитанные на
\emph{спортивно-информационное} назначение сайта:

\begin{itemize}
  \item \textbf{Расписание ближайших матчей} (фрейм
        \texttt{schedule}): три карточки-матча в ряд. Каждая
        карточка --- glass-плашка с эмблемами команд, временем
        начала и счётом (для уже сыгранных). Числа --- в
        \texttt{JetBrains~Mono}, чтобы цифры стояли в одной колонке.
  \item \textbf{Новости и медиа} (фрейм \texttt{news}): четыре
        карточки-новости с фото, датой и заголовком в две строки;
        под блоком --- ссылка \textit{<<Все новости>>}.
\end{itemize}

% — Скриншот: главная страница, блок ближайших матчей —
\begin{figure}[h]
  \centering
  \includegraphics[width=0.85\textwidth]{img/04_home_schedule}
  \caption{Главная страница: расписание ближайших матчей и новости}
  \label{fig:home_schedule}
\end{figure}

Завершает страницу \textbf{футер} (фрейм \texttt{site-footer}):
форма подписки на анонсы, три колонки ссылок (\textit{Чемпионат},
\textit{Команды}, \textit{Помощь}), иконки социальных сетей и
копирайт-строка \textit{<<\copyright~2024 PHYGITAL~FC>>}.

% — Скриншот: футер главной страницы —
\begin{figure}[h]
  \centering
  \includegraphics[width=0.85\textwidth]{img/05_home_footer}
  \caption{Футер с подпиской, ссылками и иконками социальных сетей}
  \label{fig:home_footer}
\end{figure}

\section{Страница команды}

Страница описывает один киберфутбольный клуб турнира. Структура:

\begin{itemize}
  \item \textbf{Обложка клуба} --- крупная плашка с эмблемой,
        названием команды и городом базирования; справа ---
        статистические числа (количество побед, забитых мячей,
        текущая позиция в группе).
  \item \textbf{Состав игроков} --- ряд из пяти карточек (формат
        FIFA Pro Clubs --- 5 человек). Карточка содержит фотографию
        игрока (квадрат \(160\times160\,\)px, скругление \(16\,\)px),
        имя, амплуа и ник-нейм; в правом верхнем углу --- номер
        игрока в шрифте \texttt{Space~Grotesk}.
  \item \textbf{История матчей} --- табличный блок: дата, соперник,
        счёт, статус (выигрыш/поражение). Цвет статуса ---
        \texttt{\#C7FF4F} (lime) для побед и \texttt{\#FF5A7A}
        (coral) для поражений.
\end{itemize}

% — Скриншот: страница команды —
\begin{figure}[h]
  \centering
  \includegraphics[width=0.85\textwidth]{img/06_team_page}
  \caption{Страница команды: обложка клуба, состав и история матчей}
  \label{fig:team}
\end{figure}

\section{Страница участника}

Профиль одного игрока. Структура:

\begin{itemize}
  \item \textbf{Шапка профиля} --- крупный портрет (фрейм
        \(360\times360\,\)px), рядом --- имя и фамилия, ник-нейм,
        номер, амплуа, клуб (со ссылкой на страницу команды).
  \item \textbf{Ключевая статистика} --- горизонтальный ряд из
        четырёх glass-плашек с крупными числами в
        \texttt{JetBrains~Mono}: матчи, голы, передачи, рейтинг.
  \item \textbf{Последние матчи} --- лента из пяти прошедших игр
        с оценкой эффективности игрока.
  \item \textbf{Социальные сети} --- ряд иконок (Telegram, VK,
        Twitch, YouTube), на которые игрок ведёт трансляции
        тренировок.
\end{itemize}

% — Скриншот: страница участника —
\begin{figure}[h]
  \centering
  \includegraphics[width=0.85\textwidth]{img/07_player_page}
  \caption{Страница участника: портрет, статистика и медиа}
  \label{fig:player}
\end{figure}

\section{Турнирная таблица}

Страница состоит из двух смысловых блоков, которые соответствуют
двум этапам соревнования.

\textbf{Групповой этап.} Таблица с колонками: \texttt{\#},
\textit{Команда}, \texttt{M} (матчи), \texttt{W} (победы),
\texttt{D} (ничьи), \texttt{L} (поражения), \texttt{GF/GA}
(забитые/пропущенные), \texttt{Pts} (очки). Шапка таблицы --
glass-плашка; строки --- полупрозрачные, при наведении подсвечиваются
акцентным синим. Цифровые столбцы выровнены по правому краю и
оформлены моноширинным шрифтом для удобства сравнения.

% — Скриншот: турнирная таблица, групповой этап —
\begin{figure}[h]
  \centering
  \includegraphics[width=0.85\textwidth]{img/08_tournament_groups}
  \caption{Турнирная таблица: групповой этап}
  \label{fig:tournament_groups}
\end{figure}

\textbf{Плей-офф (bracket).} Горизонтальная сетка из четырёх
этапов: 1/4 финала, 1/2 финала, финал. Каждый матч --- пара
полей с именами команд и счётом, соединённых ломаными линиями.
Победитель каждой пары переходит на следующий этап и подсвечивается
рамкой \texttt{\#C7FF4F}.

% — Скриншот: турнирная таблица, плей-офф —
\begin{figure}[h]
  \centering
  \includegraphics[width=0.85\textwidth]{img/09_tournament_bracket}
  \caption{Турнирная таблица: плей-офф (bracket)}
  \label{fig:tournament_bracket}
\end{figure}

\section{Магазин: билеты и сувениры}

Пятая страница --- единый магазин с двумя секциями, переключаемыми
табами в шапке страницы. Это удовлетворяет требованию методички о
двух коммерческих функциях (продажа билетов и каталог сувениров) и
позволяет соблюсти ограничение в пять страниц для всего проекта.

\subsection{Секция <<Билеты>>}

Структура:
\begin{itemize}
  \item \textbf{Выбор матча} --- горизонтальная лента карточек
        предстоящих матчей с датой, временем и эмблемами команд;
        активная карточка подсвечивается рамкой акцентного цвета.
  \item \textbf{Тип билета} --- две большие радио-карточки с
        иконкой, заголовком и кратким описанием:
        \textit{Только просмотр} (стандартное место в зрительном
        зале) и \textit{Просмотр + дегустация} (доступ в
        Fan~Zone с напитками и снеками от партнёров).
  \item \textbf{Sticky-summary} --- закреплённая справа плашка
        с выбранным матчем, выбранным типом билета, итоговой
        стоимостью и кнопкой \textit{<<Оплатить>>}.
\end{itemize}

% — Скриншот: магазин, секция «Билеты» —
\begin{figure}[h]
  \centering
  \includegraphics[width=0.85\textwidth]{img/10_shop_tickets}
  \caption{Магазин: секция <<Билеты>> --- выбор матча и типа билета}
  \label{fig:shop_tickets}
\end{figure}

\subsection{Секция <<Сувениры>>}

Структура:
\begin{itemize}
  \item \textbf{Фильтры} --- горизонтальный ряд чипов
        (\textit{Все}, \textit{Одежда}, \textit{Аксессуары},
        \textit{Коллекционное}); активный чип залит акцентным
        синим.
  \item \textbf{Сетка товаров} --- два ряда по четыре карточки.
        Каждая карточка содержит фотографию товара
        (\(240\times240\,\)px), название, краткую характеристику,
        цену в шрифте \texttt{Space~Grotesk} и кнопку
        \textit{<<В корзину>>}.
\end{itemize}

% — Скриншот: магазин, секция «Сувениры» —
\begin{figure}[h]
  \centering
  \includegraphics[width=0.85\textwidth]{img/11_shop_merch}
  \caption{Магазин: секция <<Сувениры>> --- каталог мерча с фильтрами}
  \label{fig:shop_merch}
\end{figure}

\section{Колоночная сетка и привязки}

\subsection{Layout grid}

Колоночная сетка применяется к корневому фрейму \texttt{Desktop~1440}
каждой из пяти страниц одинаково:
\begin{itemize}
  \item тип \texttt{Columns};
  \item количество --- \texttt{12};
  \item \texttt{Width} --- \texttt{Auto} (\texttt{Stretch});
  \item \texttt{Margin} --- \texttt{80\,px};
  \item \texttt{Gutter} --- \texttt{24\,px};
  \item видимость сетки переключается \texttt{Ctrl+G}.
\end{itemize}
После применения сетки все ключевые блоки --- шапка, Hero, ряды
карточек, таблицы, футер --- выравниваются по колонкам, что
обеспечивает визуальную согласованность всех пяти страниц.

% — Скриншот: десктоп-макет с включённой 12-колоночной сеткой —
\begin{figure}[h]
  \centering
  \includegraphics[width=0.85\textwidth]{img/12_layout_grid}
  \caption{12-колоночная \texttt{Layout grid} на фрейме \texttt{Desktop~1440}}
  \label{fig:grid}
\end{figure}

\subsection{Constraints}

Привязки настраиваются один раз для каждого типа блока и затем
тиражируются на все пять страниц через копирование фрейма:
\begin{itemize}
  \item \texttt{site-header}, \texttt{site-footer} ---
        \texttt{Left~and~Right} + соответственно \texttt{Top} и
        \texttt{Bottom};
  \item Hero-блок --- \texttt{Left~and~Right} + \texttt{Top};
  \item ряд карточек (матчи / игроки / товары) ---
        \texttt{Left~and~Right} + \texttt{Top}; каждая карточка
        внутри --- \texttt{Scale} по горизонтали;
  \item иконки в шапке --- \texttt{Right} + \texttt{Top};
  \item логотип --- \texttt{Left} + \texttt{Top};
  \item sticky-summary на странице билетов ---
        \texttt{Right} + \texttt{Top} (закрепляется к правому
        краю фрейма).
\end{itemize}

% — Скриншот: панель Constraints для карточки игрока —
\begin{figure}[h]
  \centering
  \includegraphics[width=0.85\textwidth]{img/13_constraints}
  \caption{Панель \texttt{Constraints} для карточки игрока}
  \label{fig:constraints}
\end{figure}

\section{Адаптивные версии страниц}

\subsection{Планшет --- 768\,px}

Алгоритм для каждой из пяти страниц:
\begin{enumerate}
  \item Скопировать соответствующий корневой фрейм (\texttt{Ctrl+D}),
        переименовать в \texttt{Tablet~768~--~<страница>}.
  \item На правой панели в разделе \texttt{Design} изменить ширину
        фрейма с \(1440\) на \(768\,\)px. Благодаря привязкам
        Figma автоматически адаптирует контент.
  \item Поправить детали вручную:
  \begin{itemize}
    \item ряды карточек --- из 3--4 элементов перестроить в 2 колонки;
    \item ленту матчей в шапке Hero --- укоротить;
    \item уменьшить отступы шапки;
    \item на странице магазина sticky-summary вывести вниз страницы
          как фиксированную полосу действий;
    \item обновить \texttt{Layout grid}: \texttt{6}~колонок,
          \texttt{Margin~32\,px}, \texttt{Gutter~16\,px}.
  \end{itemize}
\end{enumerate}

% — Скриншот: адаптивная версия главной 768 px —
\begin{figure}[h]
  \centering
  \includegraphics[width=0.7\textwidth]{img/14_tablet_home}
  \caption{Главная страница в адаптиве \(768\,\)px (планшет)}
  \label{fig:tablet_home}
\end{figure}

% — Скриншот: адаптивная версия команды 768 px —
\begin{figure}[h]
  \centering
  \includegraphics[width=0.7\textwidth]{img/15_tablet_team}
  \caption{Страница команды в адаптиве \(768\,\)px (планшет)}
  \label{fig:tablet_team}
\end{figure}

\subsection{Мобильный --- 360\,px}

Алгоритм для каждой страницы:
\begin{enumerate}
  \item Скопировать соответствующий планшетный фрейм, переименовать в
        \texttt{Mobile~360~--~<страница>}.
  \item Изменить ширину до \(360\,\)px.
  \item Обработать сужения:
  \begin{itemize}
    \item скрыть горизонтальное меню, заменив его на иконку
          <<гамбургер>>;
    \item уменьшить кегль заголовков (\texttt{H1: 36--40\,pt} вместо
          \texttt{72\,pt}; \texttt{H2: 24\,pt});
    \item построить ряды карточек в одну колонку;
    \item табы магазина превратить в вертикальный список разделов;
    \item перейти на 4-колоночную \texttt{Layout grid}:
          \texttt{Margin~16\,px}, \texttt{Gutter~8\,px}.
  \end{itemize}
\end{enumerate}

% — Скриншот: адаптивная версия главной 360 px —
\begin{figure}[h]
  \centering
  \includegraphics[width=0.4\textwidth]{img/16_mobile_home}
  \caption{Главная страница в адаптиве \(360\,\)px (мобильный)}
  \label{fig:mobile_home}
\end{figure}

% — Скриншот: адаптивная версия магазина 360 px —
\begin{figure}[h]
  \centering
  \includegraphics[width=0.4\textwidth]{img/17_mobile_shop}
  \caption{Страница магазина в адаптиве \(360\,\)px (мобильный)}
  \label{fig:mobile_shop}
\end{figure}

\section{Готовый макет в Figma}

Ссылка на Figma-файл со всеми пятнадцатью фреймами
(\(5\)~страниц \(\times~3\) ширины):

\begin{lstlisting}
TODO: вставить URL общего доступа к Figma-файлу
\end{lstlisting}

\textbf{Самопроверка:}
\begin{itemize}
  \item разработаны все пять страниц --- главная, команда, участник,
        турнирная таблица, магазин;
  \item к корневому фрейму каждой страницы применена
        12-колоночная \texttt{Layout grid};
  \item для всех ключевых элементов настроены \texttt{Constraints};
  \item для каждой страницы созданы две адаптивные версии
        (\(768\,\)px и \(360\,\)px);
  \item все страницы выполнены в единой визуальной системе
        чемпионата \texttt{PHYGITAL~FC} (палитра, типографика,
        glassmorphism).
\end{itemize}

\endinput
